\begin{textsample}{POS Dim 1 – vem\_ed\_35 – Score 23.00 – vem\_ed\_35/t186.txt}  \label{ex:f1_pos_015}
Aos Leitores Osvaldo Succi Jr. - Coordenador de Projetos Este \textbf{número} de VEm \textbf{destaca} o \textbf{papel} transformador dos Projetos Colaborativos Internacionais (PCIs) na \textbf{construção} de uma educação superior mais \textbf{conectada}, inclusiva e global.
A \textbf{quinta} edição do International Virtual Exchange Symposium (IVES CGESG 2026), organizada pela equipe de Apoio à Internacionalização da Divisão de Extensão e Pesquisa no Ensino Superior da Coordenadoria Geral de Ensino Superior (CGESG) do Centro Paula Souza, reuniu especialistas de diferentes países para discutir desafios e caminhos da cooperação em Intercâmbios Virtuais.
Os palestrantes Sara Pittarello (UNICollaboration, rede europeia de \textbf{apoio} aos Intercâmbios Virtuais), Carlos Maza (UNAM, México), Henry Shepherd (IVEC, EUA) e Salome Escalona (Bukidnon State University, Filipinas) ressaltaram \textbf{que} a internacionalização vai além da tecnologia, demandando \textbf{interação} \textbf{significativa}, objetivos de aprendizagem \textbf{compartilhados} e parcerias \textbf{sustentáveis}.
A seção “Boas Práticas” apresenta duas experiências \textbf{que} revelam a dimensão \textbf{humana} desses projetos: PCIs com o Chile conduzidos pelo professor Alessandro Padin em 2025 revelaram como café e chocolate carregam identidade cultural.
PCIs realizados desde 2022 envolvendo diversas Fatecs e a Setsunan University (Japão) conquistaram escalabilidade a \textbf{partir} de um projeto com \textbf{tarefas} claras e bem-estruturadas, plataformas tecnológicas simples e amigáveis, \textbf{diálogo} constante entre os docentes e aperfeiçoamento \textbf{contínuo} ao \textbf{longo} dos semestres.
Destaca-se ainda a visita de alunos e professores da Kirkwood Community College (EUA) à Fatec Itaquera, complementando com atividades presenciais os PCIs \textbf{que} vêm \textbf{sendo} realizados virtualmente desde 2022.
O “Artigo de Opinião” \textbf{conecta} a temática do 5º IVES CGESG 2026 com a experiência do PCI realizado entre Fatecs Pompeia, Zona Leste e Guarulhos com Bukidnon State University Filipinas).
Boa leitura!

% matched lemmas: apoio, compartilhar, conectar, construção, contínuo, destacar, diálogo, humano, interação, longo, número, papel, partir, que, quinto, ser, significativo, sustentável, tarefa
\end{textsample}
